\section{Certified Approximation, Locality, and Pruning}\label{sec:physics}

The formal docking story used by DQDock is not a single theorem but a chain of theorem families. The first layer supplies explicit tail and approximation bounds. The second converts those bounds into optimizer-invariance statements. The third turns optimizer invariance into locality and structural-rank control. The fourth carries those results into sampled docking and ranking-safe pruning. The importance of this layering is methodological: each family discharges one step in the argument from physical approximation to engine-level certified elimination.

\subsection{Tail Bounds and Approximation Radii}

The lattice-sum development proves explicit tail decay for Lennard--Jones-style interactions. In particular, the handle family \leanmeta{\LHrng{LS}{1}{4}.} supplies the mathematical tail substrate from which cutoff error estimates are derived. The Lennard--Jones and Coulomb families \leanmeta{\LHrng{LJ}{1}{6}, \LHrng{CB}{1}{4}.} then package exact-versus-coarse comparisons for finite action/state domains. These theorems do not claim that every physically relevant error term has already been fully discharged from first principles. Rather, they establish a clean interface: once one has a valid tail-decay or convergence bound, the decision-theoretic consequences are theorem-level.

In writing terms, this is where the paper can be most precise. Theorems \leanmeta{\LHrng{LS}{1}{2}.} give explicit asymptotic tail control; \leanmeta{\LHrng{LJ}{1}{2}, \LHrng{CB}{1}{2}.} move from that control to finite-domain uniform approximation; and \leanmeta{\LH{LJ3}, \LH{APX3}.} turn uniform approximation into exact optimizer preservation under a strict-gap condition. That chain is already enough to justify a clean formal slogan: bounded physical approximation error becomes bounded decision error, and bounded decision error becomes zero decision error once the margin is large enough.

The abstract approximation family \leanmeta{\LHrng{APX}{1}{3}.} is useful here because it isolates the logic from the chemistry. If exact and coarse scorers differ uniformly by at most $\delta$, then a finite-domain worst-case discrepancy radius exists and serves as a canonical approximation witness. Whenever the strict utility gap exceeds $2\delta$, the coarse scorer and exact scorer induce the same winner. This is the mathematical form of a statement often used informally in docking: approximation is harmless if it is small relative to the energy gap. The present framework makes that statement exact.

\subsection{From Bounded Perturbation to Locality}

The bounded-potential bridge \leanmeta{\LHrng{BP}{1}{3}.} converts physical perturbation control into a decision-locality theorem. If outside-cutoff perturbations are bounded by a tail term, and if the finite minimum strict utility gap is positive, then choosing a sufficiently large cutoff forces the perturbation error below half the minimum gap. At that point the outside-cutoff coordinates are decision-irrelevant. This is the formal step that turns physical decay into exact optimizer preservation.

That bridge matters because it upgrades a numerical cutoff heuristic into a semantic statement about the optimizer. The cutoff is no longer justified merely because it seems physically reasonable or computationally convenient. It is justified because the omitted interaction tail is too small to move the decision across the relevant gap. In that precise sense, the cutoff is a certified abstraction rather than a blind truncation.

The molecular structural-rank family \leanmeta{\LHrng{MD}{1}{7}.} then makes the docking consequence explicit. Outside-cutoff coordinates are irrelevant; therefore only atoms inside the retained pocket can contribute to the effective decision dimension. The resulting structural-rank bound is pocket-local rather than ambient. This is one of the central conceptual claims of the paper: the formal complexity of docking is controlled not by total molecular size alone, but by decision-relevant pocket structure.

The paper should lean hard on this point. It is what makes DQDock more than a certified scoring note. The engine is meant to exploit the fact that the optimizer often depends on a much smaller pocket-local support than the ambient state description would suggest. Structural rank is the formal object that measures that support, and the docking story is strongest when it is told in those terms rather than only in terms of chemistry intuition.

\subsection{Sampled Docking and Certified Top-k Guarantees}

Practical docking engines do not optimize over a literal continuous action space. They work with sampled or generated candidate poses, then re-rank and refine them. The sampled-docking families \leanmeta{\LHrng{SD}{1}{9}.} carry the exact/coarse and locality arguments into that finite candidate setting. Under the stated support and compatibility hypotheses, sampled docking inherits winner-preservation, locality, and structural-rank control from the ambient exact theory.

This transport layer is what lets the paper speak directly about an implemented docking engine rather than only about an idealized molecular decision problem. The exact ambient problem is still the semantic reference point, but the finite candidate family used in practice becomes a theorem-governed proxy for it. Put differently, the sampled-docking theorems explain when finite search is merely incomplete exploration and when it remains a faithful decision-preserving restriction.

The ranking and pruning families \leanmeta{\LHrng{TK}{1}{6}, \LH{CP1}.} then turn this into a practical engine guarantee. Uniform score error bounds imply pairwise order preservation whenever exact gaps dominate approximation error. They also imply top-$k$ containment, ambiguity-band control near the boundary, and sound pruning certificates. In operational terms, a pose may be discarded only when the theorem says it is safely outside the certified survivor set. This is the core link between the mathematical theory and the DQDock implementation.

This is also where the engine's certification language becomes most operational. A certified top-$k$ statement is not simply an abstract ranking theorem; it is exactly the kind of statement needed to justify keeping or discarding candidate poses before expensive refinement. The pruning certificate is therefore the most implementation-facing theorem family in the current stack.

\subsection{Discretization and Grid Transport}

The final bridge is discretization. DQDock uses finite representations, sampled poses, and discrete kernels even when the ambient chemistry is naturally continuous. The grid/discretization family \leanmeta{\LHrng{GD}{1}{3}.} gives the formal shape needed to connect the two. Lipschitz-style control of utility transport to a grid implies a resolution-controlled approximation theorem, which in turn implies the uniform-approximation statements needed by the ranking and winner-preservation layer. This does not solve every continuous-to-discrete issue in docking, but it gives the exact theorem template the paper can build around.

Taken together, these families justify the paper's central mathematical claim: DQDock is built around a compositional proof spine. Tail bounds support gap-safe approximation; gap-safe approximation supports locality; locality supports low structural rank; low structural rank supports sampled and discretized reductions; and those reductions support certified ranking and pruning.
